\chapter{Restriction to a plane and a line}
\label{cr21:ch}

The first higher scalar class survives restriction to a plane and a line.
Its image contains the discriminant torsion of the plane coefficient.
This torsion disappears after the plane splits into two lines.
The two restrictions therefore require different localizations.

The potentials, normalized vanishing cycles, complex orientations, and
ordinary atlas pullbacks are those of Section~\ref{sch:sec:scalar}.
All cohomology has rational coefficients and the analytic topology.
The groups below retain their full Borel coefficient complexes.
Every localization is algebraic. No completion occurs.

\section{The scalar support map}
\label{cr21:sec:support}

Let $G=\operatorname{GL}_3(\mathbb C)$ and
$L=\operatorname{GL}_2(\mathbb C)\times\mathbb C^*\subset G$.
Write $U$ for the defining plane and $Q$ for the defining line.
Their ordered direct sum gives
\[
 j:V_2\times V_1\longrightarrow V_3,
 \qquad f_3j=f_2+0.
\]
Put $K_{3,G}=R\Gamma(BG,P_{3,0})$ and
$K_{2,L}=R\Gamma(BL,P_{2,0}\otimes\mathbb Q[2])$.
The notation for $K_{2,L}$ includes the line coefficient.
Define
\[
 R=H^*(BL;\mathbb Q)=\mathbb Q[c_1,c_2,y],
 \quad c_i=c_i(U),\quad y=c_1(Q),
 \quad |c_i|=2i,\quad |y|=2.
\]
Scalar translation identifies every coincident scalar point with the origin.

For a compact subgroup $K\subset L$, write
$X\mathbin{/\mkern-6mu/}K=EK\times_KX$.
Choose the product Hilbert--Schmidt norm on all matrix entries.
It is invariant under unitary conjugation and phase rotation of $B$.
Let $F_n=f_n^{-1}(1)$ and
$P_n^{\leq}=\{\operatorname{Re}f_n\leq0\}$.
The support realization of normalized vanishing cycles gives
\begin{equation}
 \begin{aligned}
 K_{3,L}&\simeq
   \operatorname{Fib}\bigl(C^*(BL)\longrightarrow
                      C^*(F_3\mathbin{/\mkern-6mu/}L)\bigr)[18],\\
 K_{2,L}&\simeq
   \operatorname{Fib}\bigl(C^*(BL)\longrightarrow
             C^*((F_2\times V_1)\mathbin{/\mkern-6mu/}L)\bigr)[10].
 \end{aligned}
 \label{cr21:eq:fibres}
\end{equation}
Here Borel constructions for $L$ use its maximal compact subgroup.
Here $K_{3,L}=R\Gamma(BL,P_{3,0}|_L)$ restricts the equivariant
coefficient before taking $L$-equivariant cohomology.

Indeed, positive radial contraction identifies the global support pair
with every sufficiently small centered pair.
The cube root on $\{\operatorname{Re}z>0\}$ trivializes the complement
as the potential-one fibre times this half-plane.
Both maps commute with the compact group and the block inclusion.
This proves \eqref{cr21:eq:fibres} as fibre complexes, including their maps.
The support construction is Massey\cite{schMasseyST}*{\S1, p.~2}.

Restriction of these pairs, preceded by restriction from $G$ to $L$,
defines the actual degree-eight map
\begin{equation}
 \theta_{2,1}:H^k(K_{3,G})\longrightarrow H^{k+8}(K_{2,L}).
 \label{cr21:eq:theta}
\end{equation}
Its map on Milnor complements is the pullback under
$F_2\times V_1\hookrightarrow F_3$.
Equivalently, it is the scalar restriction of
\[
 j^*\mathcal P_3\longrightarrow
       (\mathcal P_2\boxtimes\mathcal P_1)[8].
\]
Vanishing-cycle pullback and the zero-function Thom--Sebastiani
identification construct this morphism before cohomology.
See Davison\cite{cr11Davison}*{\S2.2, equations (19)--(23), pp.~10--11}.
Thus the source of \eqref{cr21:eq:theta} contains the entire scalar
critical coefficient, including its equivariant extension data.

\section{The Euler denominator and the target torsion}
\label{cr21:sec:target}

Let $E_{2,1}\subset V_3\times\operatorname{Gr}(2,3)$ be the
incidence of triples preserving a plane.
Its projection $p:[E_{2,1}/G]\to[V_3/G]$ is proper.
The block projection $q:[E_{2,1}/G]\to[(V_2\times V_1)/L]$
has the split section $s$, and $ps=j$.
Its relative tangent complex is
\[
 [\,W\longrightarrow W^{\oplus3}\,],
 \qquad W=\operatorname{Hom}(Q,U),
\]
in degrees $-1,0$.
For block matrices $A_U,A_Q$, and similarly $B,C$, its differential is
\[
 t\longmapsto(A_Ut-tA_Q,B_Ut-tB_Q,C_Ut-tC_Q),
\]
up to a common sign.
The potential on the incidence is $q^*(f_2+0)$.
Smooth pullback of vanishing cycles and the contractible upper-entry
and unipotent fibres identify $q^*$ with an isomorphism on cohomology.
Its inverse is $s^*$.
This is an ambient incidence statement. The critical fibre product
retains all upper-entry commutator equations.

Set
\begin{equation}
 d=c_1^2-4c_2,
 \qquad r=y^2-c_1y+c_2,
 \qquad e=r^2.
 \label{cr21:eq:euler}
\end{equation}
If $x_1,x_2$ are the Chern roots of $U$, then
$d=(x_2-x_1)^2$ and $r=(y-x_1)(y-x_2)$.
The complex-oriented Euler ratio of the relative tangent is
$e_L(W^{\oplus3})/e_L(W)=r^2$.
It has cohomological degree eight.
The normalized scalar restriction is therefore
\begin{equation}
 \Delta_{2,1}^{\mathrm{loc}}=e^{-1}\theta_{2,1}:
 H^*(K_{3,G})\longrightarrow H^*(K_{2,L})[e^{-1}].
 \label{cr21:eq:delta}
\end{equation}
It has degree zero.
The identity $s^*p^*=(q^*)^{-1}p^*$ identifies it with the localized
Hall restriction of Jindal--Kaubrys--Latyntsev\cite{cr11JKL}*{\S3.2.3,
Lemma~3.2.4, pp.~26--27} in these ranks and conventions.

Let $v$ denote the scalar class of $\sigma_2\boxtimes\sigma_1$.
Let $t$ denote the class of $\tau_2\boxtimes\sigma_1$, where
$\tau_2$ inserts the quotient-line Chern class in the plane flag.
Put $w=t-c_1v$.

\begin{proposition}[The plane coefficient after adjoining a line]
\label{cr21:prop:target}
There is a decomposition of graded $R$-modules
\begin{equation}
 H^*(K_{2,L})=Rv\oplus(R/(d))w,
 \qquad |v|=-6,\quad |w|=-4.
 \label{cr21:eq:module}
\end{equation}
Multiplication by $e$ is injective on this module.
In particular, $w$ survives localization at $e$.
The normalized restriction of the plane to two lines sends
$v$ to $1\otimes1\otimes1$ and $w$ to zero.
\end{proposition}

\begin{proof}
First retain the diagonal torus in $\operatorname{GL}_2$.
Theorem~\ref{cr11:thm:scalar} computes its actual fibre complex and gives
one free class and one class annihilated by $(x_2-x_1)^2$.
The flag formula of Theorem~\ref{cr11:thm:hall}, with the quotient
line decorated, gives
\[
 \theta_{1,1}(t)=d(x_1+x_2),\qquad
 \theta_{1,1}(v)=d.
\]
Hence $w=t-c_1v$ belongs to the torsion summand.
Lemma~\ref{sch:lem:rank-two-decoration} identifies the ordinary image
of $t$ with the nonzero positive flag decoration $w_2$.
The ordinary image of $c_1v$ is zero.
Thus $w$ is the nonzero generator of that summand in its lowest degree.

Restriction from $\operatorname{GL}_2$ to its torus identifies rational
equivariant cohomology with Weyl invariants in the torus cohomology.
This statement also holds with the present equivariant complex:
the flag fibration and Leray--Hirsch use the tautological line class
and the same pulled-back coefficient.
The classes $v,t,w$ descend under $\operatorname{GL}_2$.
Taking invariants is exact over $\mathbb Q$.
The invariant coefficient ring is $\mathbb Q[c_1,c_2]$, and the
invariant torsion quotient is $\mathbb Q[c_1,c_2]/(d)$.
Adjoining $y$ and the line shift $[2]$ proves \eqref{cr21:eq:module}.

The free summand has no $e$-torsion.
In the other summand, $c_2=c_1^2/4$ and
\[
 e\bmod d=(y-c_1/2)^4\ne0
       \quad\text{in }R/(d)=\mathbb Q[c_1,y].
\]
This quotient is a domain, so multiplication by $e$ is injective there.
The plane restriction follows from its two displayed images.
\end{proof}

The ordinary plane-and-line coefficient has degrees $-6,-4,-1$.
The last degree comes from the equivariant extension in
\eqref{cr11:eq:borel-triangle}.
Tensoring the cohomology module \eqref{cr21:eq:module} with
$R/(c_1,c_2,y)$ loses that class.
Forgetting stabilizers therefore acts on the original derived coefficient.
There is no augmentation $R[e^{-1}]\to\mathbb Q$ extending
$c_1,c_2,y\mapsto0$.
The raw ordinary scalar pullback is zero on cohomology in every degree.
Indeed, the source vanishes below degree $-6$, whereas the shifted
target $(P_{2,0}\otimes\mathbb Q[2])[8]$ has degrees $-14,-12,-9$.
Thus polynomial evaluation of selected normalized images is a separate
operation from forgetting stabilizers in the raw support map.

\section{Restriction of the complete-flag morphism}
\label{cr21:sec:flags}

Write $\nu_n(p)$ for the complete-flag coefficient morphism with
a polynomial $p$ in the successive line Chern classes inserted.
Its scalar equivariant value has the shifts of $\nu_n$ and the
cohomological degree of $p$.
For the internal plane flag, put $s=c_1(S)$ and $u=c_1(U/S)$.
Then
\begin{equation}
 s+u=c_1,\qquad su=c_2,
 \qquad \nu_2(1)=2v,\qquad \nu_2(u)=c_1v+w.
 \label{cr21:eq:plane-flag}
\end{equation}
The notation on the right includes the external line coefficient.

\begin{lemma}[The three positions of the external line]
\label{cr21:lem:flag}
For every polynomial $p$ in three ordered line Chern classes,
\begin{equation}
 \theta_{2,1}\nu_3(p)
 =e\left\{\nu_2\bigl(p(y,s,u)\bigr)
           +\nu_2\bigl(p(s,y,u)\bigr)
           +\nu_2\bigl(p(s,u,y)\bigr)\right\}.
 \label{cr21:eq:flag}
\end{equation}
The equality holds in $H^*(K_{2,L})$, including its torsion summand.
\end{lemma}

\begin{proof}
Use the central circle in $L$ which acts trivially on $U$ and
by scalar multiplication on $Q$.
Its fixed complete flags have three components, each $\mathbb P(U)$.
They are indexed by the position of $Q$ in the complete flag.
Only the two cross-block pairs move. The internal plane flag remains.

Localization at $r$ suffices for this fixed-locus calculation.
Here is an explicit concentration argument on the flag variety.
Besides the three fixed components, the $L$-orbits have representatives
\[
 \begin{split}
 &(\langle e_1\rangle\subset\langle e_1,e_2+q\rangle),\\
 &(\langle e_1+q\rangle\subset\langle e_1,q\rangle),\\
 &(\langle e_1+q\rangle\subset\langle e_1+q,e_2\rangle),
 \end{split}
\]
where $e_1,e_2$ span $U$ and $q$ spans $Q$.
Projection to $U,Q$ and the intersection of the plane with $U$
give this list directly.
On the first stabilizer, $U/\langle e_1\rangle\simeq Q$.
On the other two, $\langle e_1\rangle\simeq Q$.
Consequently $\operatorname{Hom}(Q,U)$ has a trivial line subquotient,
so its top Chern class $r$ restricts to zero on every such stabilizer.
The finite orbit filtration and its relative-cohomology exact sequences
kill the complement after inverting $r$.
This argument applies to equivariant coefficient complexes and to
their pullbacks to the incidence.
Over a fixed flag, a nonzero moving matrix entry gives a nowhere-zero
section of $\operatorname{Hom}(Q,U)$ or its dual.
Its Euler class is again $r$ and vanishes on that locus.
The six possible cross-block maps give a finite open cover of the
remaining nonfixed incidence. Mayer--Vietoris gives the same concentration.

The complex Thom self-intersection formula now expresses the proper
incidence map as the sum of its three fixed contributions.
One can apply it directly to the support pairs defining $\Phi$:
the Thom maps, excision maps, and proper direct images commute with
restriction to $\{\operatorname{Re}f\leq0\}$.
Thus the coefficient on each fixed component is the actual $\nu_2$.
It is not the constant coefficient on the block critical locus.
This compatibility is also the construction in
Davison\cite{cr11Davison}*{\S2.7, Proposition~2.16, pp.~21--22}.

For each cross pair, let $\lambda$ be the tangent weight which
moves the earlier line towards the later line.
The three forbidden lower entries each have weight $\lambda$.
Their Thom Euler class divided by the flag tangent Euler class
is $\lambda^3/\lambda=\lambda^2$.
The two cross pairs therefore contribute $r^2=e$, independent
of the position of $Q$.
Every orientation is complex, and all Thom degrees are even.
The three restrictions of $p$ are the three displayed substitutions.
This proves \eqref{cr21:eq:flag} after inverting $r$.
Inverting $r$ and $e=r^2$ gives the same ring.
Proposition~\ref{cr21:prop:target} makes localization of the target
injective, so the equality holds before localization.
\end{proof}

Put $C=c_1+y$, the restriction of the first Chern class of the
rank-three defining representation.
The polynomial identity \eqref{cr21:eq:plane-flag} gives
\[
 \nu_2(s)=c_1v-w.
\]
Inserting $1$ and each of the three line classes in
\eqref{cr21:eq:flag} therefore gives
\begin{equation}
 \begin{array}{c|c}
 \text{class}&\Delta_{2,1}^{\mathrm{loc}}(\text{class})\\\hline
 \sigma_3&v\\
 \kappa_{3,1}&2Cv-2w\\
 \kappa_{3,2}&2Cv\\
 \kappa_{3,3}=2\tau_3&2Cv+2w.
 \end{array}
 \label{cr21:eq:values}
\end{equation}
For example, the last line is
$\nu_2(u)+\nu_2(u)+y\nu_2(1)=2Cv+2w$.
The normalization $\sigma_3=\nu_3(1)/6$ gives the first line.

\begin{theorem}[The first two scalar generators]
\label{cr21:thm:generators}
Let $z_3=\tau_3-C_1\sigma_3\in H^{-4}(K_{3,G})$, where
$C_1\in H^2(BG)$ is the first Chern class.
Then
\begin{equation}
 \Delta_{2,1}^{\mathrm{loc}}(\sigma_3)=v,
 \qquad \Delta_{2,1}^{\mathrm{loc}}(z_3)=w\ne0.
 \label{cr21:eq:generators}
\end{equation}
The restriction has polynomial image on the
$H^*(BG)$-submodule generated by $\sigma_3$ and $z_3$.
The second image has annihilator $(d)$ in $R$ and survives
cross-block localization.
\end{theorem}

\begin{proof}
Divide \eqref{cr21:eq:flag} by $e$ and use
\eqref{cr21:eq:values}.
The map is linear over $H^*(BG)$ through block restriction.
Equation~\eqref{cr21:eq:module} proves the annihilator assertion.
\end{proof}

Polynomiality in this theorem concerns the specified submodule.
Divisibility of $\theta_{2,1}(H^*(K_{3,G}))$ by $e$ in the full
module \eqref{cr21:eq:module} is a further question.
The map \eqref{cr21:eq:delta} is defined on the full source
regardless of that divisibility.

\section{Compatibility with the actual Hall image}
\label{cr21:sec:product}

The scalar incoming coefficient for a plane and a line is
\[
 D_0=\left(Rb_{2,1*}a_{2,1}^*
       (\mathcal P_2\boxtimes\mathcal P_1)\right)_0.
\]
Its $G$-equivariant version retains the tautological plane and quotient
on $\operatorname{Gr}(2,3)$.
Let $a^G$ be its lowest class multiplied by $c_1(Q)$.
Let $b^G$ use the actual morphism $\tau_2$ in the plane factor.
These are equivariant lifts of the geometric classes $a,b$ in
Theorem~\ref{sch:thm:degree-four-image}.

Refining the plane by its complete flag gives identities of coefficient
morphisms, before taking scalar cohomology:
\begin{equation}
 \mu_{2,1}(a^G)=\frac12\nu_3(\ell_3)=\tau_3,
 \qquad \mu_{2,1}(b^G)=\nu_3(\ell_2)=\kappa_{3,2}.
 \label{cr21:eq:refine}
\end{equation}
The factor $1/2$ is the normalization $\sigma_2=\nu_2(1)/2$.
Proper refinement, cup product, and the complex Thom maps commute
on the same support pairs.
Consequently \eqref{cr21:eq:values} computes the composite on these
actual incoming classes:
\begin{equation}
 \begin{aligned}
 \Delta_{2,1}^{\mathrm{loc}}\mu_{2,1}(a^G)&=Cv+w,\\
 \Delta_{2,1}^{\mathrm{loc}}\mu_{2,1}(b^G)&=2Cv.
 \end{aligned}
 \label{cr21:eq:product}
\end{equation}

There is also an exact equivariant identity behind the zero ordinary
image of $b$:
\begin{equation}
 \kappa_{3,2}=2C_1\sigma_3
          \quad\text{in }H^{-4}(K_{3,G}).
 \label{cr21:eq:middle}
\end{equation}
To prove it, use the scalar degrees
$H^{-6}(P_{3,0})=\mathbb Q$, $H^{-5}(P_{3,0})=0$, and
$H^{-4}(P_{3,0})=\mathbb Q$.
They were computed from the commuting-pair strata in
Proposition~\ref{sch:prop:scalar} and
Lemma~\ref{sch:lem:next-degrees}.
The Borel spectral sequence in total degree $-4$ has only
$H^2(BG)H^{-6}(P_{3,0})$ and $H^0(BG)H^{-4}(P_{3,0})$.
Neither supports a differential: the intervening scalar group and
odd base cohomology vanish, and lower scalar groups are zero.
Thus its dimension is two.
The images $Cv$ and $Cv+w$ of $C_1\sigma_3$ and $\tau_3$
are linearly independent over $\mathbb Q$.
The map is therefore injective in this degree.
Equation~\eqref{cr21:eq:values} proves \eqref{cr21:eq:middle}.

Forgetting stabilizers in \eqref{cr21:eq:refine} and
\eqref{cr21:eq:middle} gives
\[
 \mu_{2,1}(a)=w_3,\qquad \mu_{2,1}(b)=0.
\]
These are the actual ordinary degree-$-4$ images, with the same
positive flag normalizations as Theorem~\ref{sch:thm:degree-four-image}.
Equation~\eqref{cr21:eq:product} is their equivariant refinement.
Only its polynomial outputs can be evaluated at $c_1=c_2=y=0$.
They evaluate to $w$ and zero, respectively.
In this degree $w$ has the ordinary image $w_2\otimes v_1$.
This evaluation does not define an augmentation of
\eqref{cr21:eq:delta} on its full localized target.

The proper map for the multiplication remains
\[
 q_1:I=E_{2,1}\cap f_3^{-1}(1)\longrightarrow F_3.
\]
The varieties $I,F_3$ are smooth of complex dimensions $22,26$.
Their Gysin map has degree eight.
Together with the support boundary maps, it is the map
$H^5(I)\to H^{13}(F_3)$ of
Proposition~\ref{sch:prop:degree-four}.
The restriction uses the separate closed inclusion
$F_2\times V_1\hookrightarrow F_3$ and cohomological pullback.
Thus \eqref{cr21:eq:product} compares a proper Gysin map with a
specified pullback of the same support coefficient.
The affine projection to noncommuting pairs is unnecessary.
Its projected incidence map is nonproper, as the family following
Proposition~\ref{sch:prop:degree-four} shows.

\section{The first coassociative flag}
\label{cr21:sec:coassoc}

Let $T_3=(\mathbb C^*)^3$, with ordered Chern roots $x_1,x_2,x_3$,
and put
\[
 R_3^T=\mathbb Q[x_1,x_2,x_3],\qquad
 D=\prod_{1\leq i<j\leq3}(x_j-x_i)^2.
\]
The direct inclusion of three scalar lines defines a raw support map
\[
 \theta_{1,1,1}:H^k(K_{3,G})\longrightarrow (R_3^T)^{k+18}.
\]
It is the first arrow of the vanishing-cycle triangle after restriction
to the diagonal triples.
Define $\Delta_{1,1,1}^{\mathrm{loc}}=D^{-1}\theta_{1,1,1}$,
with target $R_3^T[D^{-1}][6]$ and degree zero.

\begin{theorem}[Three-line coassociativity]
\label{cr21:thm:coassoc}
As maps on the full scalar source $H^*(K_{3,G})$,
\begin{equation}
 (\Delta_{1,1}^{\mathrm{loc}}\otimes1)\Delta_{2,1}^{\mathrm{loc}}
 =\Delta_{1,1,1}^{\mathrm{loc}}
 =(1\otimes\Delta_{1,1}^{\mathrm{loc}})\Delta_{1,2}^{\mathrm{loc}}.
 \label{cr21:eq:coassoc}
\end{equation}
Both iterated maps take values in $R_3^T[D^{-1}][6]$.
Here $\Delta_{1,2}^{\mathrm{loc}}$ uses the reversed block order
and its corresponding Euler factor.
On $\sigma_3,\tau_3,z_3$, the common values are respectively
$1$, $x_1+x_2+x_3$, and $0$.
\end{theorem}

\begin{proof}
The two composites of block inclusions are the same ordered diagonal
inclusion $V_1^3\hookrightarrow V_3$.
Pullbacks of the support pairs compose with this identity.
The zero-function Thom--Sebastiani maps are the unit identifications,
so the two raw composites are $\theta_{1,1,1}$.
All inserted shifts are even and give no Koszul sign.
This is the scalar support realization of the coassociativity in
Jindal--Kaubrys--Latyntsev\cite{cr11JKL}*{Lemma~2.2.2, pp.~19--20}.

The two products of Euler denominators are equal:
\[
 \begin{aligned}
 &(x_2-x_1)^2
       \bigl((x_3-x_1)(x_3-x_2)\bigr)^2=D,\\
 &(x_3-x_2)^2
       \bigl((x_2-x_1)(x_3-x_1)\bigr)^2=D.
 \end{aligned}
\]
They have degree twelve, the total raw shift.
Each constituent Euler factor becomes invertible in the displayed
target ring, so both composites are defined there.
This proves \eqref{cr21:eq:coassoc}.
Proposition~\ref{cr21:prop:target} sends $v$ to $1$ and $w$ to zero.
Apply it to \eqref{cr21:eq:generators} and
$\Delta_{2,1}^{\mathrm{loc}}(\tau_3)=Cv+w$.
\end{proof}

For the complete scalar flag, let
$X=\ell_2+\ell_3$ and $Y=\ell_3$.
These are the two degree-two quotient Chern classes on
$\operatorname{Fl}(1,2;3)$.
Its actual coefficient map sends them to
$\kappa_{3,2}+\kappa_{3,3}$ and $\kappa_{3,3}$.
Their plane-and-line restrictions are
\[
 4Cv+2w,\qquad 2Cv+2w.
\]
The ordinary Hall images are both $2w_3$.
The common three-line restrictions are $4(x_1+x_2+x_3)$ and
$2(x_1+x_2+x_3)$, respectively.
This computes the first higher three-constituent flag while preserving
the equivariant terms which vanish on the ordinary atlas.

The construction establishes the scalar rank-$(2,1)$ restriction and
the stated three-line compatibility.
A full polynomial restriction requires divisibility on every scalar degree.
A comparison with holomorphic Chern--Simons observables requires a
separate morphism preserving these coefficients and correspondences.
